\documentclass[12pt]{article}
\usepackage[T1]{fontenc}
\usepackage{lmodern}
\usepackage{amsmath}
\usepackage{amssymb}
\usepackage{xcolor}
\usepackage[normalem]{ulem}
\usepackage{graphicx}
\usepackage{float}
\usepackage[autostyle]{csquotes}
\usepackage{geometry}
\usepackage{hyperref}
\geometry{a4paper, margin=1in}

\usepackage[style=apa, backend=biber]{biblatex}
\addbibresource{/Users/tio/Documents/GitHub/September_UMEDCA/references.bib}
% % Setup for bibliography
% \usepackage[style=apa, backend=biber]{biblatex}

% % Creating a bib file using the keys requested by the user, based on the provided messy bibtex file.
% % This ensures the LaTeX document is self-contained and uses the correct citation keys.
% \begin{filecontents}{reconstructed_references.bib}
% @book{inwood1992hegel,
%   author      = {Inwood, Michael James},
%   title       = {A Hegel Dictionary},
%   year        = {1992},
%   publisher   = {Blackwell Oxford}
% }

% @book{Hegel1830Logic,
%   author      = {Hegel, Georg Wilhelm Friedrich},
%   title       = {Hegel's Logic},
%   year        = {1830},
%   publisher   = {Marxists Internet Archive}
% }

% @book{Pippin:2019aa,
%   author      = {Pippin, Robert B},
%   title       = {Hegel's Realm of Shadows: Logic as Metaphysics in ``The Science of Logic''},
%   year        = {2019},
%   publisher   = {University of Chicago Press}
% }

% % Based on the source context (Carspecken (ed) Critical Ethnography and Education.pdf)
% @book{Carspecken1995aa,
%   author      = {Carspecken, Phil Francis},
%   title       = {Critical Ethnography and Education},
%   year        = {1995},
%   publisher   = {Routledge}
% }

% @book{Carspecken1999,
%   author      = {Carspecken, Phil Francis},
%   title       = {Four Scenes for Posing the Question of Meaning and Other Essays},
%   year        = {1999},
%   publisher   = {Peter Lang Inc., International Academic Publishers}
% }

% @book{Brandom2008,
%   author      = {Brandom, Robert B.},
%   title       = {Between Saying and Doing: Towards an Analytic Pragmatism},
%   year        = {2008},
%   publisher   = {Oxford University Press}
% }

% @book{Brandom2000,
%   author      = {Brandom, Robert B.},
%   title       = {Articulating Reasons: An Introduction to Inferentialism},
%   year        = {2000},
%   publisher   = {Harvard University Press}
% }

% @book{brandom1994making,
%   author      = {Brandom, Robert B.},
%   title       = {Making It Explicit: Reasoning, Representing, and Discursive Commitment},
%   year        = {1994},
%   publisher   = {Harvard University Press}
% }

% @book{brandom_reason_2013,
%   author      = {Brandom, Robert B.},
%   title       = {Reason in Philosophy: Animating Ideas},
%   year        = {2013},
%   publisher   = {Belknap Press}
% }

% @book{Brandom:2019aa,
%   author      = {Brandom, Robert B.},
%   title       = {A Spirit of Trust: A Reading of Hegel's Phenomenology},
%   year        = {2019},
%   publisher   = {The Belknap Press of Harvard University Press}
% }

% @book{Derrida1997,
%   author      = {Derrida, Jacques},
%   title       = {Of Grammatology},
%   year        = {1997},
%   translator = {Spivak, Gayatri Chakravorty},
%   publisher   = {The Johns Hopkins University Press}
% }

% @book{Derrida2011,
%   author      = {Derrida, Jacques},
%   title       = {Voice and Phenomenon},
%   year        = {2011},
%   publisher   = {Northwestern University studies in phenomenology and existential philosophy}
% }

% @book{Lucy2004Derrida,
%   author      = {Lucy, Niall},
%   title       = {A Derrida Dictionary},
%   year        = {2004},
%   publisher   = {Blackwell Publishing}
% }

% @book{agamben_language_2006,
%   author      = {Agamben, Giorgio},
%   title       = {Language and Death: The Place of Negativity},
%   year        = {2006},
%   publisher   = {University Of Minnesota Press}
% }

% @book{Derrida:1973,
%   author      = {Derrida, Jacques},
%   title       = {Speech and Phenomena},
%   year        = {1973},
%   publisher   = {Northwestern University Press}
% }

% % Note: The provided messy bib entry lists the publication year as 2010.
% @book{Hegel:2014aa,
%   author      = {Hegel, Georg Wilhelm Friedrich},
%   title       = {The Science of Logic},
%   year        = {2010},
%   publisher   = {Cambridge University Press}
% }

% @book{hegel1977phenomenology,
%   author      = {Hegel, Georg Wilhelm Friedrich},
%   title       = {Phenomenology of Spirit},
%   year        = {1977},
%   translator = {Miller, A. V.},
%   publisher   = {Oxford University Press}
% }

% @book{caygill_kant_2004,
%   author      = {Caygill, Howard},
%   title       = {A Kant dictionary},
%   year        = {2004},
%   publisher   = {Wiley-Blackwell}
% }
% \end{filecontents}

% \addbibresource{reconstructed_references.bib}

% Custom Macros from Appendix A for context
\newcommand{\SBoxUp}{\Box_{\mathsf{S}}^{\uparrow}}
\newcommand{\SBoxDown}{\Box_{\mathsf{S}}^{\downarrow}}
\newcommand{\OBoxUp}{\Box_{\mathsf{O}}^{\uparrow}}
\newcommand{\OBoxDown}{\Box_{\mathsf{O}}^{\downarrow}}
\newcommand{\SDiaUp}{\Diamond_{\mathsf{S}}^{\uparrow}}

\title{Grounding the Modal Logic of Embodiment: A Literature Review of Appendix A}
\author{Reconstructed Analysis}
\date{October 17, 2025}

\begin{document}

\maketitle

\begin{abstract}
"Appendix A: Modal Logic and Embodied Reasoning" outlines an ambitious program to provide formal, mathematical frameworks connecting phenomenology, Hegelian dialectics, and inferential semantics through a novel system of polarized modal logic. This document grounds the claims made in Appendix A within the relevant scholarly literature, demonstrating how the proposed formalizations intersect with the work of G.W.F. Hegel, Robert Brandom, Phil Carspecken, and Jacques Derrida. The analysis focuses on the reconstruction of Hegel's foundational triad through embodied phenomenology, the mapping of inferential commitments and entitlements onto embodied dynamics of compression and expansion, and the broader implications for modeling finite/infinite tension.
\end{abstract}

\section{Introduction}

The document "Appendix A: Modal Logic and Embodied Reasoning" seeks to formalize concepts developed narratively in a larger work, utilizing "polarized modal logic" to capture the directional, embodied nature of possibility and necessity. The core innovation is the polarization of standard modal operators ($\Box$ and $\Diamond$) to reflect fundamental rhythms of Temporal Compression ($\downarrow$: contraction, fixation) and Temporal Decompression ($\uparrow$: expansion, release). The appendix specifically aims to map material inferences onto these embodied dynamics; interpret the Zeeman Catastrophe Machine (ZCM) as a model of finite/\textit{in}finite tension; and reconstruct Hegel's foundational triad (Being, Nothing, Becoming) through embodied phenomenology. This review situates these formal claims within the established philosophical traditions they engage.

\section{Hegelian Dialectics and Embodied Phenomenology}

A central claim of Appendix A is the reconstruction of Hegel's foundational triad—Being, Nothing, and Becoming—through the lens of embodied phenomenology, utilizing the proposed modal logic to articulate the necessity of the dialectical movement.

\subsection{The Foundational Triad}

The sequence of Being, Nothing, and Becoming forms the starting point of Hegel's \textit{Science of Logic} (SL) \parencite{inwood1992hegel, Hegel1830Logic}. This triad describes the fundamental movement from simple, indeterminate concepts to complex, determinate ones. Understanding this logical "movement" and the "speculative logic of negation" is crucial for interpreting the work's progress \parencite[p. 38]{Pippin:2019aa}. Appendix A attempts to formalize this movement by interpreting the attempt to grasp "Being" as an act of intense Temporal Compression ($\downarrow$) and conceptual crystallization ($\OBoxDown$). The instability of this indeterminate concept leads to its necessary liquefaction ($\OBoxUp$) into "Nothing." "Becoming" emerges not as a static state, but as the movement itself—the recognition and release of this oscillation ($\SBoxUp$).

\subsection{Embodied Grounding}

The attempt to ground these abstract conceptual movements in felt, embodied experience resonates strongly with phenomenological approaches in critical theory, particularly the work of Phil Carspecken. Carspecken seeks to reconstruct philosophical concepts, including those of Hegel (specifically concerning the phylogenesis of communication), through detailed phenomenological examinations of meaningful acts \parencite{Carspecken1999}.

Carspecken argues that meaning is fundamentally tied to the body and that meaningful acts originate in "action impeti," which involve felt, bodily components. He states explicitly that "Meaning is an embodied phenomena, always with correlates in 1 st person body sensations" \parencite[p. 20]{Carspecken1995aa}. These impeti include a "proprioceptive configuration felt in the body" and an "I-feeling mode" \parencite[p. 20]{Carspecken1995aa}. By mapping Hegelian dialectics onto dynamics of tension and release, Appendix A aligns with Carspecken's insistence on the body as the correlate of holistic meaning, suggesting that the necessity of the dialectic is not merely logical but felt.

\section{Inferentialism and Polarized Modal Logic}

Appendix A proposes its polarized modal logic to formalize the dynamics of reasoning, explicitly mapping these onto the inferentialist framework developed by Robert Brandom. The core innovation is interpreting inferential roles through embodied dynamics: compression ($\downarrow$) as commitment-strengthening and expansion ($\uparrow$) as entitlement-weakening.

\subsection{Commitment and Entitlement}

Brandom's inferential semantics centers on the deontic statuses of \textit{commitment} and \textit{entitlement}. The meaning of a claim is determined by the commitments it undertakes and the entitlements it licenses \parencite{brandom1994making, Brandom2000}. Appendix A's formalization directly adopts this terminology but gives it an embodied, temporal interpretation. Making a determinate claim (undertaking a commitment) is formalized as "Temporal Compression" or "inferential crystallization." Exploring consequences or generalizations (entitlements) is formalized as "Temporal Decompression" or "inferential liquefaction."

\subsection{Modal Logic and Material Inference}

The use of modal logic to articulate the structure of reasoning is also characteristic of Brandom's project. He emphasizes "material inference"—inferences good in virtue of their content rather than their form. He defines these material relations of incompatibility and consequence as those holding "in virtue of what is expressed by non-logical vocabulary" \parencite[Note 9]{brandom_reason_2013}. Appendix A's polarized operators aim to capture how these material inferences, including material incompatibility (determinate negation) \parencite{Brandom:2019aa}, are navigated in embodied practice.

Furthermore, the appendix's goal of providing a rigorous mathematical framework mirrors Brandom's own methodology. Brandom often utilizes technical appendices to back up his philosophical claims. In the preface to \textit{Between Saying and Doing}, he describes these formal sections (which articulate incompatibility semantics for modal logic) as:
\begin{quote}
the cash for the promissory notes offered by my descriptions of those results in the lecture. As I remark there, in this context, proof is the word made flesh. \parencite[p. xix]{Brandom2008}
\end{quote}
Appendix A operates in a similar spirit, attempting to make the embodied word into formalized proof.

\section{Tension, Finitude, and Metaphysics}

The appendix also engages with the modeling of existential tension, particularly through the Zeeman Catastrophe Machine (ZCM), and touches upon themes relevant to the critique of metaphysics.

\subsection{Modeling Finite/Infinite Tension}

Appendix A interprets the ZCM as a formal model of the tension between the finite (the determinate object of attention, fixation) and the infinite (the ground of Self-Certainty). The dynamics of the ZCM—the buildup of tension (Compressive Necessity, $\SBoxDown$), the unstable cusp region (Expansive Possibility, $\SDiaUp$), and the catastrophic "snap" or release (Expansive Necessity, $\SBoxUp$)—are mapped onto the embodied experience of dissonance, choice, and sublation (the "Aha!" moment).

This modeling relates to Hegel's distinction between "bad infinity" (endless alternation without qualitative transformation, like the Being/Nothing loop) and "true infinity" \parencite{inwood1992hegel}. In Hegel's logic, the "infinite" is understood as transcending the "finite," not by annulling it, but by providing the "conceptual space within which the finite can emerge in its multifarious forms and yet also be contained by the infinite" \parencite[p. xxi]{Hegel:2014aa}. The ZCM model formalizes how the drive toward the finite necessarily generates the tension that, when released, results in a reintegration with the infinite ground.

\subsection{Derrida and the Critique of Presence}

The appendix's engagement with formalization, embodiment, and the infinite inevitably intersects with Jacques Derrida's critique of the "metaphysics of presence" \parencite{Derrida2011}. As noted in the translator's preface to \textit{Of Grammatology}, "Derrida uses the word 'metaphysics' very simply as shorthand for any science of presence" \parencite[p. xxi]{Derrida1997}. This tradition, according to Derrida, is characterized by the "determination of being as presence in all the senses of this word," prioritizing essence and fixed principles \parencite[p. xxi, quoting Derrida]{Derrida1997}.

Hegel's philosophy, particularly the \textit{Phenomenology of Spirit} \parencite{hegel1977phenomenology}, is often seen by deconstructionists as the culmination of this tradition \parencite{Derrida:1973, agamben_language_2006}. Appendix A's attempt to formalize the dialectic might be viewed critically from a Derridean perspective as an attempt to eliminate \textit{chance} and secure a stable ground within a "scientific" discourse determined by metaphysics \parencite{Lucy2004Derrida}. However, the appendix's emphasis on rhythm, movement, and the necessary instability of fixation (liquefaction) also resonates with Derrida's concept of \textit{différance}—the intertwining of spacing and deferral that keeps meaning always in movement \parencite{Derrida:1973}.

\section{Broader Connections: Space, Time, and Recognition}

Appendix A introduces a notation combining vertical polarity (temporal compression/expansion) with horizontal polarity (spatial forward derivation/retroactive recognition). This structure aims to formalize distinct philosophical insights regarding the recognition of enabling conditions. The leftward-expansive form ($A \leftarrow B$) captures the structure of Kant's transcendental deduction, where the unity of apperception is recognized retroactively through the immediacy of experience \parencite{caygill_kant_2004}. Similarly, Carspecken highlights how forestructures build through proprioceptive expansion and are recollected as enabling conditions \parencite{Carspecken1999}. Brandom's work on anaphoric chains also emphasizes the retroactive establishment of meaning \parencite{brandom1994making}. The framework thus locates embodied reasoning within a two-dimensional space of time and space, derivation and recognition.

\section{Conclusion}

"Appendix A: Modal Logic and Embodied Reasoning" proposes a novel framework that synthesizes insights from disparate philosophical traditions. It grounds the abstract movements of the Hegelian dialectic in the embodied phenomenology articulated by Carspecken, interpreting the abstract movement of the Logic as a felt necessity of compression and release. It adopts the inferentialist vocabulary of Brandom, mapping commitments and entitlements onto these same embodied dynamics. The polarized modal logic offers a precise vocabulary for analyzing reasoning not as disembodied symbol manipulation, but as an embodied, rhythmic navigation of tension between the finite and the infinite.

\printbibliography

\end{document}